\chapter*{แผนการจัดการเรียนรู้แบบมุ่งเน้นผลลัพธ์ OBE รายวิชา 4011111}
\addcontentsline{toc}{chapter}{แผนการจัดการเรียนรู้แบบมุ่งเน้นผลลัพธ์ OBE รายวิชา 4011111}

\noindent\textbf{ข้อมูลทั่วไปของรายวิชา}
\begin{itemize}
    \item \textbf{รหัสและชื่อรายวิชา}\quad 4011111 ฟิสิกส์สำหรับจุลชีววิทยา (\textbf{Physics for Microbiology})
    \item \textbf{จำนวนหน่วยกิต}\quad 3(2-3-4) หน่วยกิต (บรรยาย 2 ชั่วโมงต่อสัปดาห์ ปฏิบัติการ 3 ชั่วโมงต่อสัปดาห์ และศึกษาด้วยตนเอง 4 ชั่วโมงต่อสัปดาห์ รวมเวลาเรียนในชั้นเรียน 5 ชั่วโมงต่อสัปดาห์ ตลอด 15 สัปดาห์)
    \item \textbf{หลักสูตรและประเภทวิชา}\quad วิทยาศาสตรบัณฑิต สาขาวิชาจุลชีววิทยา (วิชาบังคับในกลุ่มวิชาแกนวิทยาศาสตร์)
    \item \textbf{ภาคการศึกษาที่เปิดสอน}\quad ภาคเรียนที่ 2 ปีการศึกษา 2569
\end{itemize}

\noindent\textbf{คำอธิบายรายวิชาภาษาไทย}

ระบบหน่วย เวกเตอร์และการเคลื่อนที่ กฎการเคลื่อนที่ของนิวตัน งานและพลังงาน ความร้อนและอุณหพลศาสตร์ ไฟฟ้าสถิตและไฟฟ้ากระแสเบื้องต้น ทัศนศาสตร์ กัมมันตภาพรังสี และปฏิบัติการที่เกี่ยวข้องกับการประยุกต์ใช้ในบริบททางจุลชีววิทยา

\noindent\textbf{คำอธิบายรายวิชาภาษาอังกฤษ}

Units system, vectors and motion, Newton's laws of motion, work and energy, heat and thermodynamics, basic electrostatics and electric current, optics, radioactivity, and laboratory practices designed to support applications in microbiological contexts.

\vspace{0.3cm}
\noindent\textbf{1. การเชื่อมโยงผลลัพธ์การเรียนรู้ของหลักสูตร (PLOs) สู่ผลลัพธ์การเรียนรู้รายวิชา (CLOs)}

เพื่อให้การจัดการเรียนการสอนสอดคล้องกับเล่มหลักสูตรวิทยาศาสตรบัณฑิต สาขาวิชาจุลชีววิทยา (หน้า 7) อย่างเคร่งครัด รายวิชา 4011111 จึงกำหนดความเชื่อมโยงระหว่าง \textbf{PLOs} และ \textbf{CLOs} ตามหลักการ \textbf{Constructive Alignment} ดังแสดงใน\mbox{ตารางที่} 1

\begin{table}[H]
\centering
\caption{การเชื่อมโยงผลลัพธ์การเรียนรู้ของหลักสูตร (PLO1--PLO6) สู่ผลลัพธ์การเรียนรู้รายวิชา 4011111 (CLO1--CLO4)}
\small
\begin{tabularx}{\textwidth}{>{\raggedright\arraybackslash\bfseries}p{2.0cm} >{\raggedright\arraybackslash}X >{\raggedright\arraybackslash\bfseries}p{2.8cm} >{\raggedright\arraybackslash}X}
\toprule
รหัส CLO & คำอธิบายผลลัพธ์การเรียนรู้รายวิชา (CLOs) & เชื่อมโยงกับ PLO & แนวทางการวัดและประเมินผล \\
\midrule
CLO1 & อธิบายความรู้พื้นฐานทางฟิสิกส์ (ระบบหน่วย เวกเตอร์ กลศาสตร์ ความร้อน ไฟฟ้า ทัศนศาสตร์ และกัมมันตภาพรังสี) ที่เป็นวิทยาศาสตร์พื้นฐานเพื่อเชื่อมโยงกับงานด้านจุลชีววิทยาได้ & PLO2 (อธิบายความรู้พื้นฐานและประยุกต์เพื่อแก้ปัญหา) & แบบทดสอบย่อย (Quiz) ใบงานทฤษฎีแบบเน้นแนวคิด และการสอบกลางภาค--ปลายภาค \\
\addlinespace[4pt]
CLO2 & ประยุกต์ใช้หลักการทางฟิสิกส์และกระบวนการทางวิทยาศาสตร์ เพื่ออธิบายกลไกการทำงานของเครื่องมือวิทยาศาสตร์ (เช่น กล้องจุลทรรศน์, Autoclave, Centrifuge) และแก้ไขปัญหาในห้องปฏิบัติการจุลชีววิทยาได้ & PLO2 และ PLO3 (ใช้กระบวนการทางวิทยาศาสตร์แก้ปัญหาจุลชีววิทยา) & กรณีศึกษา (Case Study) การวิเคราะห์ข้อมูลการทดลอง และโครงงานย่อยสัปดาห์ที่ 14 \\
\addlinespace[4pt]
CLO3 & ปฏิบัติการทดลองทางฟิสิกส์โดยตระหนักถึงความปลอดภัย มีวินัย และปฏิบัติตามจรรยาบรรณในการทำงานในห้องปฏิบัติการ & PLO1 (มีจรรยาบรรณแห่งวิชาชีพวิทยาศาสตร์) & แบบสังเกตพฤติกรรมความปลอดภัย วินัยการใช้อุปกรณ์ และการบันทึกข้อมูลตามจริงไม่เมคเลข \\
\addlinespace[4pt]
CLO4 & ทำงานร่วมกับผู้อื่นได้อย่างมีประสิทธิภาพ และสามารถสื่อสารข้อมูลทางวิทยาศาสตร์ (การเขียนรายงานและการนำเสนอ) ได้อย่างถูกต้อง & PLO4, PLO5 และ PLO6 (พัฒนาตนเอง ทำงานเป็นทีม และสื่อสารถูกต้อง) & รายงานปฏิบัติการรายกลุ่ม การอภิปรายผล และการนำเสนอโครงงานย่อย \\
\bottomrule
\end{tabularx}
\end{table}

\vspace{0.2cm}
\noindent\textbf{2. แผนการสอนรายสัปดาห์ตลอด 15 สัปดาห์ (ทฤษฎี 2 ชั่วโมง และปฏิบัติการ 3 ชั่วโมง รวม 5 ชั่วโมงต่อสัปดาห์)}

\begin{table}[H]
\centering
\caption{แผนการจัดการเรียนรู้รายสัปดาห์ รายวิชา 4011111 ฟิสิกส์สำหรับจุลชีววิทยา (15 สัปดาห์)}
\footnotesize\setlength{\tabcolsep}{4pt}
\begin{tabularx}{\textwidth}{c >{\raggedright\arraybackslash}p{3.5cm} >{\raggedright\arraybackslash}p{3.6cm} >{\raggedright\arraybackslash}X c}
\toprule
\textbf{สัปดาห์} & \textbf{หัวข้อบรรยาย (2 ชม.)} & \textbf{หัวข้อปฏิบัติการ (3 ชม.)} & \textbf{กลยุทธ์การสอน OBE} & \textbf{CLOs} \\
\midrule
1 & บทที่ 1 ระบบหน่วยและการวัดในโลกจุลทรรศน์ ($\mu\text{m}$, $\text{nm}$, $\mu\text{L}$) & Lab 1 การวัดละเอียดและทักษะการใช้ไมโครปิเปตและการชั่งสาร & Interactive Lecture \& Hands-on Practice & CLO1, CLO3, CLO4 \\
\addlinespace[3pt]
2 & บทที่ 2 เวกเตอร์และการเคลื่อนที่ในของเหลวและการแพร่ & Lab 2 การเคลื่อนที่แนวตรงและสังเกต Brownian Motion จุลินทรีย์ & Problem-Based Learning (PBL) & CLO1, CLO3, CLO4 \\
\addlinespace[3pt]
3 & บทที่ 3 กฎการเคลื่อนที่ของนิวตัน แรงต้าน และความหนืดอาหารเหลว & Lab 3 กฎของนิวตันและการวัดความหนืดของสารละลายชีวภาพ & สาธิตและสืบเสาะ (Guided Inquiry) & CLO1, CLO2, CLO3 \\
\addlinespace[3pt]
4 & บทที่ 4 งาน พลังงาน และหลักการเครื่องปั่นเหวี่ยง (RPM สู่ RCF) & Lab 4 พลังงานกลและการคำนวณแรงเหวี่ยงสัมพัทธ์เครื่อง Centrifuge & Case-Based Learning & CLO1, CLO2, CLO4 \\
\addlinespace[3pt]
5 & บทที่ 5 ความร้อน ความจุความร้อนจำเพาะ และการควบคุมอุณหภูมิ & Lab 5 การถ่ายโอนความร้อนและตรวจสอบเสถียรภาพตู้ Incubator & Active Learning \& Group Discussion & CLO1, CLO2, CLO3 \\
\addlinespace[3pt]
6 & บทที่ 6 อุณหพลศาสตร์ ความดันไอ และหลักการหม้อนึ่งความดันไอ & Lab 6 ความดันไอ จุดเดือด และวิเคราะห์รอบการทำงาน Autoclave & Case Study (ทำไมต้อง $121\,^\circ\text{C}$ 15 psi) & CLO1, CLO2, CLO4 \\
\addlinespace[3pt]
7 & \multicolumn{3}{>{\raggedright\arraybackslash}p{9.6cm}}{\textbf{สัปดาห์ทบทวนสรุปสูตรฉบับเข้าใจง่าย (1 ชม.) และสอบวัดผลกลางภาค เนื้อหาสัปดาห์ที่ 1--6 (2 ชม.)}} & CLO1, CLO2 \\
\addlinespace[3pt]
8 & บทที่ 7 ไฟฟ้าสถิต ประจุผิวเซลล์แบคทีเรีย และการย้อมสีแกรม & Lab 7 ไฟฟ้าสถิตและกลไกการยึดเกาะของสีย้อมประจุบวกกับผนังเซลล์ & Visual Simulation \& Concept Mapping & CLO1, CLO3, CLO4 \\
\addlinespace[3pt]
9 & บทที่ 8 ไฟฟ้ากระแส กฎของโอห์ม และอิเล็กโทรโฟรีซิส & Lab 8 การต่อวงจรไฟฟ้าพื้นฐานและการจำลอง Gel Electrophoresis & Inquiry-Based Lab & CLO1, CLO2, CLO3 \\
\addlinespace[3pt]
10 & บทที่ 9 ทัศนศาสตร์พื้นฐาน การสะท้อน การหักเห และดัชนีหักเห & Lab 9 การสะท้อนและการหักเหของแสงผ่านตัวกลางและสารละลายน้ำตาล & Hands-on Optical Ray Tracing & CLO1, CLO3, CLO4 \\
\addlinespace[3pt]
11 & บทที่ 10 กล้องจุลทรรศน์ ขีดจำกัดการแยกชัด และน้ำมันช่วยส่อง & Lab 10 เลนส์นูน กล้องจุลทรรศน์ และบทบาทของ Immersion Oil (100x) & Problem-Based Lab (แก้ปัญหาภาพมัว) & CLO1, CLO2, CLO3 \\
\addlinespace[3pt]
12 & บทที่ 11 กัมมันตภาพรังสี ครึ่งชีวิต และรังสี UV ฆ่าเชื้อจุลินทรีย์ & Lab 11 การจำลองครึ่งชีวิตกัมมันตภาพรังสีและประสิทธิภาพรังสี UV-C & Simulation \& Data Plotting & CLO1, CLO2, CLO4 \\
\addlinespace[3pt]
13 & บทที่ 12 การบูรณาการฟิสิกส์กับเครื่องสเปกโตรโฟโตมิเตอร์ & Lab 12 การวัดค่าการดูดกลืนแสง ($\text{OD}_{600}$) ประเมินการเจริญแบคทีเรีย & Integrated Instrument Lab & CLO1, CLO2, CLO4 \\
\addlinespace[3pt]
14 & การสังเคราะห์ฟิสิกส์ในห้องปฏิบัติการจุลชีววิทยาสมัยใหม่ & Lab 13 โครงงานย่อย (Mini-Project) นำเสนอฟิสิกส์ในเครื่องมือจุลชีววิทยา & Project-Based Learning (PjBL) & CLO2, CLO3, CLO4 \\
\addlinespace[3pt]
15 & \multicolumn{3}{>{\raggedright\arraybackslash}p{9.6cm}}{\textbf{สัปดาห์สรุปองค์ความรู้รวบยอดและสอบวัดผลปลายภาค เนื้อหาสัปดาห์ที่ 8--14 (3 ชม.)}} & CLO1, CLO2 \\
\bottomrule
\end{tabularx}
\end{table}

\noindent\textbf{3. สัดส่วนการวัดและประเมินผลตามแนวทาง OBE (รวม 100 คะแนน)}
\begin{itemize}
    \item \textbf{ภาคทฤษฎีและการประยุกต์แนวคิด (50 คะแนน)}\quad แบ่งเป็น สอบกลางภาค 20 คะแนน (CLO1, CLO2), สอบปลายภาค 20 คะแนน (CLO1, CLO2) และแบบทดสอบย่อยท้ายคาบควบคู่ใบงานแก้โจทย์ฉบับเข้าใจง่าย 10 คะแนน (CLO1) โดยมีตารางสูตรฟิสิกส์ให้ในห้องสอบเพื่อเน้นวัดความเข้าใจและการนำไปใช้จริง
    \item \textbf{ภาคปฏิบัติการและทักษะกระบวนการทางวิทยาศาสตร์ (50 คะแนน)}\quad แบ่งเป็น รายงานปฏิบัติการรายสัปดาห์ 20 คะแนน (CLO4), ทักษะการใช้อุปกรณ์ วินัย ความปลอดภัย และการบันทึกข้อมูลตามจริง 15 คะแนน (CLO3) และโครงงานย่อยนำเสนอหลักการฟิสิกส์ในเครื่องมือจุลชีววิทยา 15 คะแนน (CLO2, CLO4)
\end{itemize}
